%! TEX program = xelatex
%! TEX root = lecture3_slide.tex

\usepackage{listings}
\title{\texorpdfstring{$\lambda$}{λ}-Calculus}
\subtitle{Parametric Polymorphism}
\begin{document}

{\usebackgroundtemplate{\includegraphics[width=\paperwidth]{banner.pdf}}
\begin{frame}\maketitle\end{frame}}

%\begin{frame}[fragile]{Polymorphism}
%  \emph{Abstraction principle~\cite{Pierce2002}}:
%  \begin{quotation}
%    Each significant piece of functionality in a program should be implemented
%    in just one place in the source code. Where similar functions are carried
%    out by distinct pieces of code, it is generally beneficial to combine them
%    into one by abstracting out the varying parts.
%  \end{quotation}
%
%  In Haskell
%  \begin{semiverbatim}
%    head :: [a] -> a
%    head (x:_) = x
%  \end{semiverbatim}
%\end{frame}

%\section{G\"odel's \textbf{T}: Simply typed $\lambda$-calculus with naturals
%and primitive recursion}
%\begin{frame}{The un defineability of $\lambda_\to$}
%  Can you write this in $\lambda_\to$ using Church numerals?
%  \begin{align*}
%    \mathit{sum}(0)     & = 0 \\
%    \mathit{sum}(1 + n) & = (1 + n) + \mathit{sum}(n)
%  \end{align*}
%  It is not definable in $\lambda_\to$, since fixpoint operator is not allowed any more. 
%
%  But, this is definable via \emph{primitive recursion}:
%  for some $c$ and function $g$
%  \begin{align*}
%    \mathit{prec}(0, c, g(x, y))       &= c \\
%    \mathit{prec}(1 + n, c, g(x, y))   &= g(n, \mathit{prec}(n, c, g(x, y)))
%  \end{align*}
%
%  $\lambda_\to$ with primitive recursion is called G\"odel's \textbf{T}.
%\end{frame}
%
%\begin{frame}{\textbf{T}: Types and terms}
%  \begin{definition}[Types]
%    \begin{multicols}{2}
%      \begin{prooftree}
%        \AXC{$B \in \mathbb{V}$}
%        \RightLabel{(tvar)}
%        \UIC{$B : \type$}
%      \end{prooftree}
%      \begin{prooftree}
%        \AXC{$\vphantom{B}$}
%        \RightLabel{(nat)}
%        \UIC{$\mathbb{N} : \type$}
%      \end{prooftree}
%      \begin{prooftree}
%        \AXC{$\sigma : \type$}
%        \AXC{$\tau   : \type$}
%        \RightLabel{(fun)}
%        \BIC{$\sigma \to \tau : \type$}
%      \end{prooftree}
%    \end{multicols}
%  \end{definition}
%  \begin{definition}[Terms]
%    Additional term formation rules are added to $\lambda_\to$ as follows.
%    \begin{multicols}{2}
%      \begin{prooftree}
%        \AXC{$\vphantom{M}$}
%        \UIC{$\zero : \term_{\textbf{T}}$}
%      \end{prooftree}
%      \begin{prooftree}
%        \AXC{$M$}
%        \UIC{$\suc\;M : \term_{\textbf{T}}$}
%      \end{prooftree}
%    \end{multicols}
%      \begin{prooftree}
%        \AXC{$L : \term_{\textbf{T}}$}
%        \AXC{$M : \term_{\textbf{T}}$}
%        \AXC{$N : \term_{\textbf{T}}$}
%        \AXC{$x \in V$}
%        \AXC{$y \in V$}
%        \QuinaryInfC{$\mathtt{rec}(M; x.\,y.\,N)\; L : \term_{\textbf{T}}$}
%      \end{prooftree}
%  \end{definition}
%\end{frame}
%
%\begin{frame}{\textbf{T}: Typing rules}
%  \begin{definition}
%    Additional term typing rules are added to $\lambda_\to$ as follows.
%    \begin{multicols}{2}
%      \begin{prooftree}
%        \AXC{$\vphantom{\Gamma}$}
%        \UIC{$\Gamma \vdash \zero : \mathbb{N}$}
%      \end{prooftree}
%      \begin{prooftree}
%        \AXC{$\Gamma \vdash M : \mathbb{N}$}
%        \UIC{$\Gamma \vdash \suc\;M : \mathbb{N}$}
%      \end{prooftree}
%    \end{multicols}
%    \begin{prooftree}
%      \AXC{$\Gamma \vdash L : \mathbb{N}$}
%      \AXC{$\Gamma \vdash M : \tau$}
%      \AXC{$\Gamma, x : \mathbb{N}, y : \tau \vdash N : \tau$}
%      \TrinaryInfC{$\Gamma \vdash \mathtt{rec}(M; x.\,y.\,N)\;L : \tau$}
%    \end{prooftree}
%  \end{definition}
%  \begin{itemize}
%    \item Substitution for \textbf{T} is defined similarly.
%    \item Substitution respects typing judgements, i.e.\ 
%      $\Gamma \vdash N : \tau$ and $\Gamma, x : \tau \vdash M : \sigma$, then 
%      $\Gamma \vdash M\subst{N}{x} : \sigma$.
%  \end{itemize}
%\end{frame}
%
%\begin{frame}{\textbf{T}: Dynamics}
%  $\beta$-conversion for \textbf{T} is extended with two rules
%  \begin{align*}
%    \mathtt{rec}(M, x.\,y.\, N)\;\zero & \longrightarrow_\beta M \\
%    \mathtt{rec}(M, x.\,y.\, N)\;\suc\;L & \longrightarrow_\beta N\subst{L , \mathtt{rec}(M; x.\,y.\,N)\;L}{x, y}
%  \end{align*}
%  Similarly, a $\beta$-reduction $\onereduce$ extends $\longrightarrow_\beta$ to all parts of a term and $\reduce$ indicates finitely many $\beta$-reductions.
%  
%  \mode<presentation>{\vfill}
%  \begin{theorem}
%    \textbf{T} enjoys the strong and weak normalisation properties as well as type safety.
%  \end{theorem}
%\end{frame}
%
%\begin{frame}{Example: Addition and summation}
%  $\mathtt{add} : \mathbb{N} \to \mathbb{N} \to \mathbb{N}$ can be defined in \textbf{T} as
%  \[
%    \lambda n.\,\lambda m.\,\mathtt{rec}\;(m; x.\,y.\, \suc\;y)\;n\;m
%  \]
%
%  $\mathtt{sum} : \mathbb{N} \to \mathbb{N}$ can be defined in \textbf{T} as
%  \[
%    \lambda n.\, \mathtt{rec}\;(\zero; x.\,y.\, \add\;(\suc\;x)\;y)\;n
%  \]
%\mode<presentation>{\vfill}
%  \begin{block}{Exercise}
%    Evaluate $\mathtt{sum}\;(\suc\;\zero)$.
%  \end{block}
%\end{frame}
%

%\section{\PCF --- System of Recursive Functions}
%
%\begin{frame}{\textbf{PCF}: $\lambda_\to$ with naturals and general recursion}
%%  \textbf{T} does not include all computable functions, since all terms terminate eventually. 
%%  Programming language in reality allows us to do \emph{general recursion}
%%  including \emph{infinite loops}. 
%
%%  \textbf{PCF} has the same class of types as \textbf{T}.
%
%  \begin{definition}[Terms]
%    Additional term formation rules are added to $\lambda_\to$ as follows.
%    \begin{multicols}{2}
%      \begin{prooftree}
%        \AXC{$\vphantom{M}$}
%        \UIC{$\zero : \term_{\PCF}$}
%      \end{prooftree}
%      \begin{prooftree}
%        \AXC{$M : \term_{\PCF}$}
%        \UIC{$\suc\;M : \term_{\PCF}$}
%      \end{prooftree}
%    \end{multicols}
%      \begin{prooftree}
%      \color{red}
%        \AXC{$L : \term_{\PCF}$}
%        \AXC{$M : \term_{\PCF}$}
%        \AXC{$N : \term_{\PCF}$}
%        \AXC{$x \in V$}
%        \QuaternaryInfC{$\ifz(M; x.\,N)\; L$}
%      \end{prooftree}
%      \begin{prooftree}
%      \color{red}
%        \AXC{$M : \term_{\PCF}$}
%        \AXC{$x \in V$}
%        \BIC{$\fix\,x.\,M : \term_{\PCF}$}
%      \end{prooftree}
%  \end{definition}
%\end{frame}
%
%\begin{frame}{\textbf{PCF}: Typing rules}
%  \begin{definition}
%    Additional term typing rules are added to $\lambda_\to$ as follows.
%    \begin{multicols}{2}
%      \begin{prooftree}
%        \AXC{$\vphantom{\Gamma}$}
%        \UIC{$\Gamma \vdash \zero : \mathbb{N}$}
%      \end{prooftree}
%      \begin{prooftree}
%        \AXC{$\Gamma \vdash M : \mathbb{N}$}
%        \UIC{$\Gamma \vdash \suc\;M : \mathbb{N}$}
%      \end{prooftree}
%    \end{multicols}
%    \begin{prooftree}
%      \color{red}
%      \AXC{$\Gamma \vdash L : \mathbb{N}$}
%      \AXC{$\Gamma \vdash M : \tau$}
%      \AXC{$\Gamma, x : \mathbb{N} \vdash N : \tau$}
%      \TrinaryInfC{$\Gamma \vdash \ifz(M; x.\,N)\;L : \tau$}
%    \end{prooftree}
%    \begin{prooftree}
%      \color{red}
%      \AXC{$\Gamma, x : \tau \vdash M : \tau$}
%      \UIC{$\Gamma \vdash \fix\,x.\,M : \tau$}
%    \end{prooftree}
%  \end{definition}
%  \begin{itemize}
%    \item Substitution for \textbf{PCF} is defined similarly.
%    \item Substitution respects typing judgements, i.e.\ 
%      $\Gamma \vdash N : \tau$ and $\Gamma, x : \tau \vdash M : \sigma$, then 
%      $\Gamma \vdash M\subst{N}{x} : \sigma$.
%  \end{itemize}
%\end{frame}
%
%\begin{frame}{\textbf{PCF}: Dynamics}
%  $\beta$-conversion for \textbf{PCF} is extended with three rules
%  \begin{align*}
%    \fix\,x.\,M & \longrightarrow_\beta M\subst{\fix\,x.\,M}{x} \\
%    \ifz(M; x.\, N)\;\zero & \longrightarrow_\beta M \\
%    \ifz(M; x.\, N)\;(\suc\;L) & \longrightarrow_\beta N\subst{L}{x}
%  \end{align*}
%  Similarly, a $\beta$-reduction $\onereduce$ extends $\longrightarrow_\beta$ to all parts of a term and $\reduce$ indicates finitely many $\beta$-reductions.
%  
%  \mode<presentation>{\vfill}
%  \begin{theorem}
%    \textbf{PCF} enjoys type safety. 
%  \end{theorem}
%\end{frame}
%
%\begin{frame}{Example: Predecessor and negation}
%  \begin{align*}
%    \pred & \defeq \lambda n : \mathbb{N}.\,\ifz(\zero; x.\, x)\;n & : \mathbb{N} \to \mathbb{N} \\
%               \mathtt{not} & \defeq \lambda n : \mathbb{N}.\, \ifz(\suc\;\zero; x.\, \zero)\; n & : \mathbb{N} \to \mathbb{N}
%  \end{align*}
%  \begin{block}{Exercise}
%    Evaluate the following terms to their normal forms.
%    \begin{enumerate}
%      \item $\pred\;\zero$
%      \item $\pred\;(\suc\;(\suc\;(\suc\;\zero)))$
%      \item $\mathtt{not}\;(\suc\;(\suc\;\zero))$
%    \end{enumerate}
%    
%  \end{block}
%\end{frame}

%\section{Polymorphic Typed \texorpdfstring{$\lambda$}{λ}-Calculus: Introduction}
%
%\begin{frame}{Polymorphism}
%
%  For example, the type $\forall t.\, t \to t$ corresponds to a type variable
%  universally quantified in GHC with extension \texttt{RankNType}
%  \begin{semiverbatim}
%    forall a. a -> a
%  \end{semiverbatim}
%  or simply 
%  \begin{semiverbatim}
%    a -> a
%  \end{semiverbatim}
%
%%\framebreak
%%GHC with the extension \texttt{ScopedTypeVariables} allows you to use type
%%variables. 
%%\begin{semiverbatim}
%% f :: forall a. [a] -> [a]
%%
%% f = \\(xs :: [a]) -> reverse xs
%%\end{semiverbatim}
%%GHC with the extension \texttt{TypeApplications} allows you to apply a type
%%argument.
%%\begin{semiverbatim}
%%  Prelude> :t id
%%
%%  id :: a -> a
%%
%%  Prelude> :t id @Int
%%
%%  id @Int :: Int -> Int
%%\end{semiverbatim}
%%
%%\framebreak
%%  Can you see the corresponding untyped/simply typed $\lambda$-terms of the
%%  following $\lambda$-terms in System F? 
%%
%%  \begin{enumerate}
%%    \item $\Lambda t.\,\lambda (x : t).\, x$
%%    \item $\Lambda t.\, \lambda (x : t)(y : t).\, x$
%%    \item $(\Lambda t.\, \lambda (x : t)(y : t).\, x)\;\sigma\;M\;N$
%%     where  $M, N$ are terms and $\sigma$ a type.
%%    \item $\Lambda t.\, \lambda (f : t \to t)(x : t).\, f\;(f\;x)$
%%  \end{enumerate}
%%System $F$ is more expressive than simply typed lambda calculus.
%\end{frame}

\section{Polymorphic \texorpdfstring{$\lambda$}{λ}-Calculus: Static}
\begin{frame}{Polymorphic types}
    Given a set $\mathbb{V}$ of type variables, the judgement $A : \type$ is defined by
    defined by
      \begin{prooftree}
        \AXC{}
        \RightLabel{(tvar), if $X \in \mathbb{V}$}
        \UIC{$X : \type$}
      \end{prooftree}
      \begin{prooftree}
        \AXC{$A : \type$}
        \AXC{$B : \type$}
        \RightLabel{(fun)}
        \BIC{$A \to B : \type$}
      \end{prooftree}
      \begin{prooftree}
        \AXC{$A : \type$}
        \AXC{$X   \in \mathbb{V}$}
        \RightLabel{(universal)}
        \BIC{$\forall X.\, A : \type$}
      \end{prooftree}
  where $X$ may or may not occur in~$A$.

  The polymorphic type $\forall X.\, A$ provides a universal type for every type $B$ by instantiating $X$ for $B$, i.e.\ $A\subst{B}{x}$.
\end{frame}

\begin{frame}{Examples}
  For example, the polymorphic type allows us to express terms that should work on arbitrary types, such as
  \begin{itemize}
    \item $\mathtt{id} : \forall X.\, X \to X$

    \item $\mathtt{proj}_1 : \forall X.\,\forall Y.\, X \to Y \to X$

    \item $\mathtt{proj}_2 : \forall X.\,\forall Y.\, X \to Y \to Y$

    \item $\mathtt{length} : \forall X.\, \List\;X \to \nat$

    \item $\mathtt{singleton} : \forall X. X \to \List(X)$
  \end{itemize}
  
\end{frame}

\begin{frame}{Free and bound variables, again}
\begin{definition}
  The \emph{free variable} $\FV(A)$ of $A$ is defined inductively by
  \begin{align*}
    \FV(X)              & = X \\
    \FV(A \to B)        & = \FV(A) \cup \FV(B) \\
    \FV(\forall X.\, A) & = \FV(A) - \{X\}
  \end{align*}
  For convenience, the function extends to contexts:
  \[
    \FV(\Gamma) = \set{ X \in \mathbb{V}}{ \exists (x : A) \in \Gamma
      \wedge X \in \FV(A) }.
  \]
\end{definition}
\begin{exercise*}

  \begin{enumerate}
    \item $\FV(\forall X.\, (X \to X) \to X \to X)$
    \item $\FV(x : X_1, y : X_2, z : \forall X.\,X)$
  \end{enumerate}
\end{exercise*}
\end{frame}

\begin{frame}{Capture-avoiding substitution for type}
  Permutation of type variables and $\alpha$-equivalence between types are defined similarly. 

  In particular, the substitution is also defined to avoid any capture of free type variables:
  \begin{definition}
  The \emph{capture-avoiding substitution} of a type $A$ for a type variable~$X$ is defined on types by 
  \begin{align*}
    X\subst{A}{X} & = A \\
    Y\subst{A}{X} & = Y && \text{if $X \neq Y$}\\
    (B \to C)\subst{A}{X} & = (B\subst{A}{X}) \to (C\subst{A}{X}) \\
    (\forall Y.B)\subst{A}{X} & = \forall Y. B\subst{A}{X} &&
    \text{if $Y \neq X, Y \not\in \FV(X)$} 
  \end{align*}
  \end{definition}
\end{frame}

\begin{frame}{Typed terms}
  Terms in polymorphic $\lambda$-calculus are extended with types.
  We define the set of terms from scratch here:
\begin{definition}
  The set $\Lambda_{\forall}(V, \mathbb{V})$ of terms in polymorphic $\lambda$-calculus is defined inductively: 
  \begin{description}
    \item[variable] $x \in \Lambda_{\forall}(V, \mathbb{V})$ if $x$ is in $V$
    \item[application] $t@u \in \Lambda_{\forall}(V, \mathbb{V})$ if $t, u \in \Lambda_{\forall}(V, \mathbb{V})$
    \item[abstraction] $\lam{x:A}t$ if $x \in V$, $A$ is a type, and $t \in \Lambda_{\forall}(V, \mathbb{V})$
    \item[type abstraction] $\lam{X}t$ is in $\Lambda_{\forall}(V, \mathbb{V})$ if $X$ is in $\mathbb{V}$ and $t$ is in $\Lambda_{\forall}(V, \mathbb{V})$
    \item[type application] $t\;A$ is in $\Lambda_{\forall}(V, \mathbb{V})$ if $t$ is in $\Lambda_{\forall}(V, \mathbb{V})$ and $A$ is a type.
  \end{description}
\end{definition}

  N.B.\ $\lam{x:A}t$ includes the type of $x$ as part of term.
  We have additionally a \emph{substitution $t\subst{A}{X}$ of a type $A$ for a type variable $X$} in $t$.
\end{frame}

\begin{frame}{Typing judgement: Overview}
  Polymorphic $\lambda$-calculus has two kinds of typing judgements.
  \begin{itemize}
    \item $\Delta \vdash A$ stands for a type $A$ under the type context $\Delta$;
    \item $\Delta; \Gamma \vdash t : A$ stands for a term $t$ of type $A$ under the context $\Gamma$ and the type context~$\Delta$
  \end{itemize}
  where a \emph{type context} is a sequence of type variable $X_1, X_2, \dots, X_n$.

  The new context $\Delta$ is used to keep track of type variables available within the term, as they may be introduced by type abstraction.

\end{frame}

\begin{frame}{Type formation}
  The judgement $\Delta \vdash A$ is constructed inductively by following rules.
  \begin{multicols}{2}
    \begin{prooftree}
      \AXC{}
      \RightLabel{if $\Delta \ni X$}
      \UIC{$\Delta \vdash X$}
    \end{prooftree}
    \columnbreak
    \begin{prooftree}
      \AXC{$\Delta \vdash X$}
      \AXC{$\Delta \vdash Y$}
      \BIC{$\Delta \vdash X \to Y$}
    \end{prooftree}
  \end{multicols}
  \begin{prooftree}
    \AXC{$\Delta, X \vdash A$}
    \UIC{$\Delta \vdash \forall X.\, A$}
  \end{prooftree}

  \begin{block}{Exercise}
    Derive the judgement 
    \[
        X \vdash X \to X
    \]
  \end{block}
\end{frame}

\begin{frame}{Typing rules}
  The judgement $\Delta; \Gamma \vdash t : A$ is defined inductively by following rules.
  \begin{multicols}{2} 
    \small
  \begin{prooftree}
    \AXC{}
    \RightLabel{if $\Gamma \ni x : A$}
    \UIC{$\Delta ; \Gamma \vdash x : A$}
  \end{prooftree}
  \begin{prooftree}
    \AXC{$\Delta; \Gamma \vdash t : A \to B$}
    \AXC{$\Delta; \Gamma \vdash u : A$}
    \BIC{$\Delta; \Gamma \vdash t\;u : B$}
  \end{prooftree}
  \begin{prooftree}
    \AXC{{\color{red}$\Delta \vdash A$}}
    \AXC{$\Delta; \Gamma, x : A \vdash t : B$}
    \BIC{$\Delta; \Gamma \vdash \lam{x : A} t : A \to B$}
  \end{prooftree}
  \color{red}
  \begin{prooftree}
    \AXC{$\Delta, X; \Gamma \vdash t : A$}
    \RightLabel{($\forall$-intro)}
    \UIC{$\Delta; \Gamma \vdash \lam{X}t : \forall X.\, A$}
  \end{prooftree}
  \begin{prooftree}
    \AXC{$\Delta; \Gamma \vdash t : \forall X.\, A$}
    \AXC{$\Delta \vdash B$}
    \RightLabel{($\forall$-elim)}
    \BIC{$\Delta; \Gamma \vdash t\;B : A\subst{B}{x}$}
  \end{prooftree}
  \end{multicols}

  \begin{theorem}[Type safety]
    Suppose $\Delta; \Gamma \vdash t : A$. Then, 
    \begin{enumerate}
      \item $t \onereduce u$ implies $\Delta; \Gamma \vdash u : A$; 
      \item $t$ is in normal form or there exists $u$ such that $t \onereduce u$
    \end{enumerate}
  \end{theorem}
\end{frame}

\begin{frame}{Undecidability of type inference}
  \begin{theorem}[Wells, 1999]
     It is undecidable whether, given a closed term $t$ of the untyped
     $lambda$-calculus, there is a well-typed term $t'$ in polymorphic $\lambda$-calculus such that
     $|t'| = t$.  
  \end{theorem}

  Two ways to retain decidable type inference:
  \begin{enumerate}
    \item Limit the expressiveness so that type inference remains decidable.
      For example, \emph{Hindley-Milner type system} adapted by Haskell 98, Standard ML, etc.\ supports only a limited form of polymorphism but type inference is decidable.
    \item Adopt \emph{partial} type inference so that type annotations can be used for, e.g. top-level definitions and 
      local definitions. 
  \end{enumerate}

  Check out \emph{bidirectional type synthesis}.
%  \begin{description}
%    \item[Arbitrary Rank Polymorphism] $\forall$ can appear
%      anywhere {\small (GHC with \texttt{-XRankNType})}. 
%    \item[Rank-1 Polymorphism]
%      $\forall$ only appear in the outermost position.
%  \end{description}
\end{frame}

\begin{frame}{Typing derivation}

The typing judgement ${}\vdash \lam{X} \lam{x : X} x : \forall X.\,X \to X$ is derivable
\vfill
\begin{prooftree}
  \LARGE 
  \AXC{}
  \UIC{$X \vdash X$}
  \AXC{}
  \UIC{$X; x : X \vdash x : X $}
  \BIC{$X ; \cdot \vdash \lam{x : X} x : X \to X$}
  \UIC{$\vdash \lam{X}\lam{x : X} x : \forall X. X \to X$}
\end{prooftree}
\vfill

\begin{convention}
  $\vdash t : A$ stands for $\cdot ; \cdot \vdash t : \tau$ where both contexts are empty.
\end{convention}
\end{frame}

\begin{frame}{Exercise}
  Derive the following judgements:
  \begin{enumerate}
    \item $\vdash (\lam{X\,Y}\lam{x : X}{y : Y} x) : \forall X.\,\forall Y.\, X \to Y \to X$
    \item $\vdash\lam{X}\lam{f : t \to t}{x : t} f\;(f\;x) : \forall X.\;(X\to X) \to X\to X$
    \end{enumerate}
  Hint.\, polymorphic $\lambda$-calculus F is syntax-directed, so the type inversion holds. 
  
\end{frame}
%\begin{frame}{Self application}
%Self-application is not typable in simply typed $\lambda$-calculus. 
%  \[
%    \lambda (x : t).\, x\; x
%  \]
%  However, self-application is possible in System F. 
%  \[
%    \lambda (x : \forall t. t\to t).\, x\;(\forall t. t\to t)\;x
%  \]
%\mode<presentation>{\vfill}
%  \begin{block}{Exercise}
%    Instantiate the first $t$ with the type $\forall t.\, t \to t$.  
%  \end{block}
%\end{frame}

\section{Polymorphic \texorpdfstring{$\lambda$}{λ}-Calculus: Dynamics and Programming}

\begin{frame}{Reduction}
  
$\beta$-reduction for polymorphic $\lambda$-calculus has two rules apart from other structural rules:
\[
  (\lam{x : A} t)\,u \onereduce t\subst{u}{x}
  \quad\text{and}\quad
  \color{red} (\lam{X}t)\;A \onereduce t \subst{A}{X}
\]

For example, 
\[
  (\lam{X}\lam{x : X} x)\;A\; t
  \onereduce
  (\lam{x : X} x)\subst{A}{X}\;t
  \equiv 
  (\lambda x : A.\, x)\;t
  \onereduce
  t
\]

Similarly, $\beta$-reduction extends to subterms of a given term, introducing relations $\onereduce$ and $\reduce$ in the same way.
\end{frame}

%\begin{frame}{System F: Evaluation}
%  
%The full $\beta$-reduction is a relation on $\lambda$-terms defined by
%\begin{multicols}{2}
%    \begin{prooftree}
%      \AXC{$M_1 \longrightarrow_{\beta} M_2$}
%      \UIC{$M_1 \onereduce M_2$}
%    \end{prooftree}
%    \begin{prooftree}
%      \AXC{$M_1 \onereduce M_2$}
%      \UIC{$\lambda (x:\tau).\, M_1 \onereduce \lambda (x:\tau).\, M_2$}
%    \end{prooftree}
%    \begin{prooftree}
%      \AXC{{\color{red}$M_1 \onereduce M_2$}}
%      \UIC{{\color{red}$\Lambda t.\, M_1 \onereduce \Lambda t.\, M_2$}}
%    \end{prooftree}
%    \begin{prooftree}
%      \AXC{$M_1 \onereduce M_2$}
%      \UIC{$M_1\,N \onereduce M_2\,N$}
%    \end{prooftree}
%    \begin{prooftree}
%      \AXC{$N_1\onereduce  N_2$}
%      \UIC{$M\,N_1 \onereduce M\,N_2$}
%    \end{prooftree}
%    \begin{prooftree}
%      \AXC{{\color{red}$M_1 \onereduce M_2$}}
%      \UIC{{\color{red}$M_1\,\tau \onereduce M_2\,\tau$}}
%    \end{prooftree}
%\end{multicols}
%If $M \onereduce N$, then $M$ and $N$ \alert{denotes} the same value. So, 
%\[
%  M =_\beta N
%\]
%where $=_\beta$ is the congruence and equivalence closure of
%$\longrightarrow_\beta$.
%\end{frame}


\begin{frame}{Sum type}

\begin{definition}
  The \emph{sum type} is defined by
  \[
    A + B \defeq \forall X. (A \to X) \to (B \to X) \to X
  \]
\end{definition}
It has two injection functions: the first injection is defined by
\begin{align*}
  \mathtt{left}_{A + B} & \defeq \lam{x : A}\lam{X}\lam{f : A \to X}{g : B\to X} f\;x \\
  \mathtt{right}_{A + B} & \defeq \lam{y : B}\lam{X}\lam{f : A \to X}{g : B \to X} g\;y
\end{align*}

%With the sum type, we can implement the usual construct \texttt{either}:
%\[
%  \mathtt{either} : \forall\,t.\;(\sigma \to t) \to (\tau \to t)
%  \to (\sigma + \tau) \to t
%\]
%which corresponds to the Haskell function
%{\small
%\begin{semiverbatim}
%  either :: (a -> c) -> (b -> c) -> Either a b -> c
%
%  either f g (Left  a) = f a 
%
%  either f g (Right b) = g b
%\end{semiverbatim}}

\mode<presentation>{\vfill}
\begin{block}{Exercise}
  Define 
  \[
    \mathtt{either} : \forall X.\, (A \to X) \to (B \to X) \to A + B \to X
  \] 
\end{block}
\end{frame}

\begin{frame}{Product type}
\begin{definition}[Product Type]
  The product type is defined by
  \[
    A \times B \defeq \forall X.\, (A \to B \to X) \to X
  \]
\end{definition}
The pairing function is defined by
\begin{align*}
  \left< \_, \_\right>_{A, B} \defeq \lam{x : A}{y : B}\lam{X}\lam{f : A \to B \to X} f\;x\;y
\end{align*}
\begin{block}{Exercise}
Define projections 
\[
  \mathtt{proj}_1 : A \times B \to A
  \quad\text{and}\quad
  \mathtt{proj}_2 : A\times B \to B
\]
\end{block}
\end{frame}
\begin{frame}[allowframebreaks]{Natural numbers}
The type of Church numerals is defined by 
\[
  \nat \defeq \forall X.\, (X\to X)\to X\to X
\]
  \begin{description}
    \item[Church numerals]
      \begin{align*}
        \bc_n & : \nat \\
        \bc_n & \defeq \lam{X}\lam{f:X \to X}{x:X} f^n\;x
      \end{align*}
    \item[Successor]
      \begin{align*}
        \suc & : \nat \to \nat \\
        \suc & \defeq \lam{n : \nat}{X}{f : X \to X}{x:X}f\;(n\;X\;f\;x) 
      \end{align*}
    \item[Addition]
      \begin{align*}
        \add & : \nat \to \nat \to \nat \\
        \add & \defeq \lam{n : \nat}{m:\nat} \lam{X}{f:X\to X}{x: X}\\
        & (m\;X\;f)\;(n\;X\;f\;x) 
      \end{align*}
    \item[Multiplication] 
      \begin{align*}
       \mul & : \nat\to \nat \to\nat\\
       \mul & \defeq\,?
      \end{align*}
    \item[Conditional]
      \begin{align*}
       \ifz & : \forall X.\,\nat \to X \to X \to X \\
       \ifz & \defeq\,?
      \end{align*}
  \end{description}
Polymorphic $\lambda$-calculus allows us to define \emph{recursor} like \texttt{fold} in Haskell.
\begin{align*}
  \mathtt{fold}_{\nat} & : \forall X.\, (X \to X) \to X \to \nat \to X  \\
  \mathtt{fold}_{\nat} & \defeq \lam{X}{f : t \to t}{e_0 : t}{n : \nat} n\;X\;f\;e_0
\end{align*}

\begin{block}{Exercise}
  Define $\add$ and $\mul$ using $\mathtt{fold}_{\nat}$ and justify your answer.
  
  \begin{enumerate}
    \item $\add'\defeq \;? : \nat\to\nat\to\nat$
    \item $\mul'\defeq \;? : \nat\to\nat\to\nat$
  \end{enumerate}
\end{block}
\end{frame}

\begin{frame}{Lists}
\begin{definition}
  For any type $A$, the type of lists over $A$ is 
\[
  \List(A) \defeq \forall X.\, X \to (A \to X \to X) \to X
\]
\end{definition}
with list constructors:
\[
  \mathtt{nil}_A \defeq \lam{X}{h : t}{f : A \to t \to t} h
\]
and $\mathtt{cons}_A$ of type $A \to \List(A) \to \List(A)$ defined as
\[
  \lam{x : A}{xs : \List(A)}{X}{h : X}{f : A \to X \to X}f\,x\,(xs\; X\; h\; f)
\]


\end{frame}

%\begin{frame}[allowframebreaks]{Existential Type}
%  Abstraction can be described by an \emph{existential type}.
%  \[
%    \exists t.\, \tau
%  \]
%  where $\tau$ is a type in which $t$ may appear as a free variable.
%  A term $(\sigma, M)$ of type $\exists t.\, \tau$ consists of a \emph{type}
%  $\sigma$ and a term $M$ of type $\tau\subst{t}{\sigma}$. 
%
%
%  \begin{example}
%    A \emph{stack} is a data structure with two operations:
%    \begin{itemize}
%      \item \texttt{push}
%      \item \texttt{pop}
%    \end{itemize}
%    Any implementation satisfying $\mathtt{pop} \circ (\mathtt{push}\, n\, st) =
%    (n, st)$ can be seen as a stack. 
%  \end{example}
%
%  \begin{example}
%    The type of stacks over natural numbers can be defined by
%    \[
%      \exists t. (t \times \nat \to t) \times (t \to 1 + \nat \times t) 
%    \]
%    \begin{enumerate}
%      \item The first component is the \text{push} function.
%      \item The second component is the \text{pop} function.
%    \end{enumerate}
%    An instance of a stack is an element $(\sigma, M)$ of the above type
%    \begin{enumerate}
%      \item $\sigma$ is a type
%      \item $\mathtt{proj}_1\, M \colon \sigma \times \nat \to \sigma$
%      \item $\mathtt{proj}_2\, M \colon \sigma \to 1 + \nat \times
%          \sigma$. 
%    \end{enumerate}
%  \end{example}
%\end{frame}
%
%%\begin{frame}[fragile]{Encoding of Existential Type in Haskell}
%%  A limited form of existential type can be defined by 
%%  \begin{semiverbatim}
%%    data Exists t where
%%      Ex :: t a -> Exists t
%%  \end{semiverbatim}
%%  and another version for typeclass
%%  \begin{semiverbatim}
%%    data Exists' c where
%%      Ex' :: c a => a -> Exists' c
%%  \end{semiverbatim}
%%  in GHC with extensions \texttt{GADTs} and \texttt{ConstraintKinds}. 
%%
%%\mode<presentation>{\vfill}
%%  \begin{block}{Exercise}
%%    Define the top type $\top$ which contains every typable term.
%%  \end{block}
%%\end{frame}
%
%\begin{frame}{Encoding of Existential Type in System F}
%  The existential type can also be encoded in System F as 
%  \[
%    \exists t.\, \tau \defeq \forall u.\; (\forall t.\; \tau \to u) \to u
%  \]
%%  using continuation passing style.\footnote{Recall that $\exists x.\,\varphi
%%  = (\forall x.\, (\varphi \to \bot)) \to \bot$ classically.  }
%\mode<presentation>{\vfill}
%  A witeness $(\sigma, M)$ of $\exists t.\, \tau$ is encoded as
%  \[
%    \Lambda u.\,\lambda (f : \forall t.\, \tau \to u).\, f\;\sigma\;M
%  \]
%
%  \begin{block}{Exercise}
%    Check that the above terms do make sense. How to use a term of some existential
%    type? 
%  \end{block}
%  
%\end{frame}

\begin{frame}{Impredicative encodings of inductive types} 
  \emph{All inductive types} can be defined in Polymorphic $\lambda$-calculus and we have demonstrated this expressiveness for sums, natural numbers, and lists, using the universal type.

\end{frame}

\begin{frame}{Type erasure}
\begin{definition}
  The \emph{erasing map} is a function defined by
  \begin{align*}
    |x| & = x \\
    |\lam{x : A} t| & = \lam{x} |t| \\
    |t\;u| & = |t|\;|u| \\
    |\lam{X} t| & = |t| \\
    |t\;A| & = |t|
  \end{align*}
\end{definition}

\begin{proposition}
  Within polymorphic $\lambda$-calculus, if ${}\vdash t : A$ and $|t| \onereduce u'$, then there exists a well-typed term~$u$ with
  ${}\vdash u : A$ and $|u| = u'$.
\end{proposition}
Well-typed terms can be evaluated with types erased safely.
\end{frame}

%\section{Parametricity}
%
%\begin{frame}{Parametricity}
%  What functions can you write for the following type?
%  \[
%    \forall X.\,X \to X 
%  \]
%  Since $X$ is arbitrary, we cannot inspect the content of $X$.
%  What we can do with $X$ is simply return it.
%  \begin{theorem}
%    Every term $t$ of type $\forall X.\, X \to X$ is \emph{observationally equivalent}%
%    \footnote{The notion of observational equivalence is beyond the scope of this lecture.}
%      to $\lam{X}{x : t} x$. 
%  \end{theorem}
%\end{frame}

%\begin{frame}{Parametricity: Theorems for free\footnote{Philip Wadler. 1989. Theorems for free! In \emph{Proceedings of the fourth international conference on Functional programming languages and computer architecture (FPCA ’89)}. ACM, New York, NY, USA, 347–359.}}
%  Assume \textbf{F} extended with the list type~$\List\;\tau$ for $\tau$ and
%  the type $\mathbb{N}$ of naturals, denoted $\mathbf{F}_{\List, \mathbb{N}}$.
%
%  Then $\mathtt{head} \circ \mathtt{map}\; f = f \circ \mathtt{head}$
%      for any $f : \tau \to \sigma$ where $\mathtt{head} : \forall t.\, \List\;t \to t$
%  can be proved by just reading the type of $\mathtt{head}$ and $\mathtt{tail}$!
%\end{frame}



%\section{Nameless Representation}
%
%\begin{frame}{Capture-avoiding but ill-defined substitution}
%  The definition of capture-avoiding substitution is not well-defined.
%  Recall that
%    \begin{align*}
%      x\subst{L}{x} & = L \\
%      y\subst{L}{x} & = y && \text{if $x \neq y$} \\
%      (M\, N)\subst{L}{x} & = M\subst{L}{x}\; N\subst{L}{x} \\
%      (\lambda x.\, M)\subst{L}{x} & = \lambda x.\, M \\
%      (\lambda y.\, M)\subst{L}{x} & = \lambda y.\, M\subst{L}{x}                                 &&\text{if $x \neq y$ and $y \not\in \FV(L)$}
%    \end{align*}
%    The function $\_\subst{L}{x}\colon \term_V \to \term_V$ is not total, so it is \alert{not} an instance of \emph{structural recursion} (i.e.\ $\texttt{fold}$).
%    In what sense, is the above well-defined? 
%    
%  \begin{enumerate}
%    \item Use \emph{nominal technique} and the notion of $\alpha$-structure recursion/induction.
%      It requires some elements of group theory.
%    \item Use \emph{nameless} representation.
%  \end{enumerate}
%  
%\end{frame}
%
%\begin{frame}[allowframebreaks]{Well-Scoped de Bruijn index representation}
%  An index $i$ starting from $0$ is used as a variable to represent the $i$-th enclosing $\lambda$ (binder) `from the inside out'.
%  For example, a term with named variables
%  \[
%    \lambda a.\,\lambda b.\, (\lambda c.\, c) \;(\lambda c.\, a\;b)
%  \]
%  becomes
%  \[
%  {\color{blue}\lambda}\, {\color{red}\lambda}\, (\lambda\,0) \;(\lambda\,{\color{blue} 2}\;{\color{red}1})
%  \]
%  Hint.\ It may be easier to think of a term in its tree representation.
%  \begin{definition}[de Bruijn representation with a local scope]
%    The term formation \fbox{$t\quad \term_n$} is defined inductively for $n \in \mathbb{N}$ by
%    \begin{minipage}{.45\textwidth}
%      \begin{prooftree}
%        \AXC{$0 \leq i < n$}
%        \UIC{$i \quad \term_{n}$}
%      \end{prooftree}
%    \end{minipage}
%    \begin{minipage}{.45\textwidth}
%      \begin{prooftree}
%        \AXC{$t \quad \term_{n + 1}$}
%        \UIC{$\lambda\, t \quad \term_n$}
%      \end{prooftree}
%    \end{minipage}
%      \begin{prooftree}
%        \AXC{$t \quad \term_{n}$}
%        \AXC{$u \quad \term_{n}$}
%        \BIC{$t\;u \quad \term_{n}$}
%      \end{prooftree}
%    \fbox{$t\quad \term_n$} means $t$ has at most $n$ many free variables.
%  \end{definition}
%  
%\end{frame}
%
%\begin{frame}{Exercise}
%  Translate the following terms to its de Bruijn index representation. 
%  \begin{enumerate}
%    \item $\lambda x.\, x$
%    \item $\lambda s.\,\lambda z.\, s \;z$
%    \item $\lambda a.\,\lambda b.\, a\;(\lambda c.\, a\; b)$
%    \item $(\lambda x.\, x)\;(\lambda y.\, y)$
%    \item $\lambda x.\, y$
%    \item $x\;y\;z$
%  \end{enumerate}
%\end{frame}
%
%\begin{frame}{Substitution, revisited}
%  How to reformulate $\beta$-reduction for terms in de Bruijn representation? 
%  Consider
%  \[
%    \left({\color{blue}\lambda}\, {\color{red}\lambda}\,(\lambda\,0) \;(\lambda\,{\color{blue} 2}\;{\color{red}1})\right)\;t
%  \longrightarrow_{\beta}
%  \left({\color{red}\lambda}\, \;(\lambda\,0) \;(\lambda\,{\color{blue} 2}\;{\color{red}1})\right)\subst{t}{\color{blue}0}
%  \]
%
%  The de Bruijn index increments under a binder so $\subst{t}{i}$ should be $\subst{t'}{i+1}$ where $t'$ is the result of incrementing every index in $t$, e.g., 
%  \begin{align*}
%      \left({\color{red}\lambda}\, (\lambda\,0) \;(\lambda\,{\color{blue} 2}\;{\color{red}1})\right)\subst{t}{\color{blue}0}
%      & = 
%      {\color{red}\lambda}\, (\lambda\,0)\subst{t'}{\color{blue}1} \quad(\lambda\,{\color{blue} 2}\;{\color{red}1})\subst{t'}{1} \\
%      & =
%      {\color{red}\lambda}\, (\lambda\,0\subst{t''}{\color{blue}2}) \quad(\lambda\,\left({\color{blue} 2}\;{\color{red}1}\right)\subst{t''}{2}) \\
%      & =
%      {\color{red}\lambda}\, (\lambda\,0) \quad(\lambda\, {\color{blue} 2}\subst{t''}{2}\;{\color{red}1}\subst{t''}{2}) \\
%      & = 
%      {\color{red}\lambda}\, (\lambda\,0) \quad(\lambda\, t''\;{\color{red}1})
%  \end{align*}
%  
%\end{frame}
%
%
%\begin{frame}{Simultaneous variable renaming}
%  \begin{definition}
%    A \emph{(variable) renaming} is a function $\rho$ between $\mathbb{Z}_n$ and $\mathbb{Z}_m$. 
%  \end{definition}
%  Every renaming $\rho\colon \mathbb{Z}_n \to \mathbb{Z}_m$ extends to an action on terms:
%  \begin{align*}
%    \left< \rho \right> i                 & = \rho(i)                                  \\
%    \left< \rho \right> \left(t\;u\right) & = \left<\rho\right>t \; \left<\rho\right>u \\
%    \left< \rho \right> \lambda\,t        & = \lambda\left< \rho' \right> t
%  \end{align*}
%  where $\rho' \colon \mathbb{Z}_{n+1} \to \mathbb{Z}_{m+1}$ is defined as 
%  \begin{align*}
%    \rho'(0) & = 0 \\
%    \rho'(1 + i) & = 1 + \rho(i)
%  \end{align*}
%  to avoid changing bound variables. 
%  
%  In particular, $\mathit{wk}\colon \term_{n} \to \term_{n+1}$ derived by $i \mapsto i+1 \in \mathbb{Z}_{n+1}$ increments every index of a free variable by $1$.
%\end{frame}
%
%\begin{frame}{Simultaneous substitution}
%  \begin{definition}
%    A \emph{(simultaneous) substitution} is a function $\sigma$ from $\mathbb{Z}_n$ to $\term_m$.
%  \end{definition}
%  Every substitution extends to an action terms: 
%  \begin{align*}
%    \left< \sigma \right> i                 & = \sigma(i)                                \\
%    \left< \sigma \right> \left(t\;u\right) & = \left<\sigma\right>t \; \left<\sigma\right>u \\
%    \left< \sigma \right> \lambda\,t        & = \lambda\left< \sigma' \right> t
%  \end{align*}
%  where $\sigma' \colon \mathbb{Z}_{n+1} \to \term_{m+1}$ is defined as
%  \begin{align*}
%    \sigma'(0) & = 0 \\
%    \sigma'(1 + i) & = \mathit{wk}\left(\sigma(i)\right)
%  \end{align*}
%\end{frame}
%
%\begin{frame}{Single substitution}
%  \begin{definition}
%    A \emph{single substitution} for $t$ is a simultaneous substitution given by
%    \begin{align*}
%      \sigma      \colon  \mathbb{Z}_{1+n} & \to \mathbb{Z}_{n} \\
%      \sigma(0)   & = t  \\
%      \sigma(1+i) & = i
%    \end{align*}
%  \end{definition}
%\end{frame}
%
%\begin{frame}{Exercise}
%  \begin{enumerate}
%    \item Adopt $\alpha$-equivalence to the de Bruijn representation.
%    \item Adopt $\beta$-equivalence to the de Bruijn representation.
%    \item Apply the new definition of substitution to compute $\mathtt{not}\;\mathtt{True}$.
%    \item Adopt the definitions of renaming and substitution to the de Bruijn level representation.
%  N.B.\ we may also count the $i$-th enclosing binder `from the outside in' using the same definition, called \emph{the de Bruijn level}.
%  \end{enumerate}
%\end{frame}

%\begin{frame}
%  $M \sim N : \mathcal{R}_{\sigma, \tau}$ is defined as follows.
%
%  \begin{enumerate}
%    \item $M \sim M : \mathcal{R}_{\mathbb{N}, \mathbb{N}}$ for every $M :
%      \mathbb{N}$.
%
%    \item For any $\mathcal{R}_{\tau, \sigma}$, 
%      \[
%        M_i \sim N_i : \mathcal{R}_{\sigma, \tau} \; \text{for any $M_i : \sigma$ and $N_i : \tau$}
%        \iff 
%        [ M_1, \dots, M_k ] \sim [ N_1, \dots, N_k ] : \mathcal{R}_{\List\;\sigma, \List\;\tau}
%      \]
%
%    \item For any $\mathcal{R}_{\sigma, \sigma'}$ and $\mathcal{S}_{\tau, \tau'}$,
%      $M \sim M' : \mathcal{R}_{\sigma, \sigma'} \implies L\; M
%      \sim L'\;M' : \mathcal{S}_{\tau, \tau'}$ for any $M$ and $M'$
%      if and only if 
%      $L \sim L' : \mathcal{R}\to\mathcal{S}_{\sigma \to \tau, \sigma' \to \tau'}$
%
%    \item Assume $\mathcal{F}$ which sends $\mathcal{R}_{\sigma, \tau}$ to
%      $\mathcal{F}\mathcal{R}_{F\sigma, F\tau}$. Then
%      $f\;\sigma \sim g\;\tau : \mathcal{F}\mathcal{R}$ for all $\mathcal{R}_{\sigma, \tau}$
%      if and only if $f \sim g : \forall \mathcal{X}_{t, u}.\, {\mathcal{F}\mathcal{X}}_{Ft, Fu}$
%  \end{enumerate}
%\end{frame}

%\begin{frame}{System $F_\omega$}
%  Recall that for $\sigma, \tau \in \type$ the sum type of $\sigma$ and $\tau$
%  is 
%  \[
%    \sigma + \tau \defeq \forall t. (\sigma \to t) \to (\tau \to t) \to t
%  \]
%  Can we internalise this \alert{type construction}? 
%  \vfill
%  \alert{\emph{Kinds}} are like \emph{classes} in set theory:
%  \begin{enumerate}
%    \item $*$ : the kind of types 
%    \item $* \Rightarrow *$ : the kind of type operators, e.g., 
%      $a \mapsto [a]$ 
%    \item \ldots 
%  \end{enumerate}
%  System $F_\omega$ is an extension of System $F$ with type-level functions and
%  kinds. See \cite{Pierce2002} for further detail. 
%\end{frame}
%
%\begin{frame}[allowframebreaks]{Impredicativity}
%  A definition is \alert{\emph{impredicative}} if it has a quantifier whose
%  domain includes itself (which is being defined). 
%  \begin{block}{Russell's Paradox}
%    \[
%      R \defeq \set{x}{ x \not\in x }
%    \]
%  \end{block}
%  Recall the self-application ...
%  \[
%    (\Lambda t.\,\lambda (x: t). x)\;{\color{red}(\forall t.\,t \to t)}\;
%    (\Lambda t.\, \lambda (x: t). x) 
%  \]
%  which is actually impredicative! 
%  This form of polymorphism is called \alert{\emph{impredicative polymorphism}}. 
%  \begin{block}{Girard's Paradox}
%    An encoding of Russell's paradox in extensions of
%    System F. 
%  \end{block}
%
%  Martin-L\"of's type theory (on which Agda is based) was inconsistent, as it
%  included an axiom\footnote{It is only a simplified story...}:
%  \[
%    \mathsf{Set} : \mathsf{Set}
%  \]
%
%  \begin{block}{Inconsistency of Haskell}
%    Impredicativity $+$ Injectivity $+$ Type Case Analysis
%    $=$ Inconsistency (Russell paradox)
%    \footnote{\url{http://okmij.org/ftp/Haskell/impredicativity-bites.html}}
%    
%  \end{block}
%\end{frame}
%\section{Parametricity}
%Polymorphism in System F is \emph{parametric} in the sense that every
%instance is uniformly defined. To put it differently, it cannot depend on any
%specific type. It is different from the \emph{ad hoc} polymorphism adopted in
%programming languages like C++ where programmers are allowed to give a
%specific implementation for some type.
%
%Why is a uniform definition important? Not only we can assure ourselves that a
%polymorphic function has a uniform result for each instance, but also it
%automatically satisfies certain properties which only depends on its type (not
%its instance).
%
%In the end of this lecture, we state a powerful result informally, originally
%called \emph{Abstraction Theorem} by Reynolds~\cite{Reynolds1983} and nowadays
%known as \emph{parametricity}. After that, we apply this theorem to characterise
%some data types introduced so far.
%
%\begin{proposition}
%  Let $\sigma$ be a type without any free variable and $\sigma^+$ the relation
%  on~$\sigma$ lifted from~$\sigma$ (defined below). Then, 
%  \[
%    \vdash F : \sigma
%    \implies (F, F) \in \sigma^+.
%  \]
%\end{proposition}
%
%To decode this proposition, we introduce a new notation: a relation $R_\sigma
%\subset \sigma_1 \times \sigma_2$ is denoted by $R_\sigma\colon \sigma_1
%\not\to\sigma_2$. For each type in System F, we define its lifted relation
%$\sigma^+$
%\begin{definition}
%  For each type formation (construct), we define its corresponding formation of
%  relations 
%  \begin{enumerate}
%    \item For $f_i:\sigma_i \to \tau_i$, $(f_1, f_2) \in R_\sigma \to S_\tau$ if
%      and only if for every $(M_1, M_2) \in R_\Sigma$ we have $(f_1\;M_1,
%      f_2\;M_2) \in S_\tau$.
%  \end{enumerate}
%\end{definition}
%
%
%
%\subsection*{Exercise}

\begin{frame}{Homework}
  \begin{enumerate}
    \item (2.5\%) Define $\mathtt{length}_\sigma : \List\;\sigma \to \nat$ calculating the length of a list in polymorphic $\lambda$-calculus.
    
  \end{enumerate}
  
\end{frame}

%\begin{frame}{References}
%\bibliographystyle{amsalpha}
%\bibliography{library} 
%\end{frame}

\end{document}
